% !TEX TS-program = lualatex
На основании проведённого обзора литературы можно сделать вывод, что применение искусственного интеллекта в контексте игр привело к значительным достижениям в этой области.

Искусственный интеллект используются для анализа огромных объемов данных. Это способствует разработке передовых стратегий и тактик, ранее не применявшихся на практике. Кит Вулси, вице-чемпион мира по нардам 1996 года, отмечал, что лишь в некоторых технических аспектах может претендовать на определённое преимущество перед искусственным интеллектом, сыгравшим миллионы партий с самим собой \cite{kitwoolsey}. Машина не подвержена эмоциональной предвзятости и учитывает только текущее состояние доски, а также потенциальные будущие результаты. Похожим образом обучаются автопилоты для самолетов и автомобилей, которые склонны минимизировать потери в случае аварии. Моральный аспект вопроса стал объектом для изучения в научных работах \cite{moralaspects}.

Искусственный интеллект может быть использован для тренировки специалистов. Программы могут анализировать игры, выявлять ошибки и предоставлять обратную связь по альтернативным ходам, способствуя обучению игроков на различных уровнях квалификации. Такие подходы также используются на производствах, для исправления бракованных продуктов в случае, если это возможно.

Исследования в области искусственного интеллекта для игр является актуальной задачей. Многие решения прошлого десятилетия могут быть улучшены благодаря вычислительным мощностям и открытиям современности. Например, некоторые представленные в настоящей работе решения, применимы только в контексте коротких нард с жестким набором правил, когда как другие имеют широкие границы применимости.

Современные подходы к использованию искусственного интеллекта в играх разнообразны и зависят от типа игры и целей исследования. Основные направления включают обучение с подкреплением, вероятностное моделирование и гибридные методы. Например, в играх с полной информацией эффективны алгоритмы, подобные AlphaZero, тогда как в играх с элементом случайности, таких как нарды, важна оценка риска и адаптивность.

В следующей главе настоящей работы будут рассмотрены ключевые подходы, применимые к решению поставленной задачи, а также проанализированы их сильные и слабые стороны в контексте выбранной предметной области.
